\documentclass[11pt]{article}

\usepackage[margin=1in]{geometry}
\usepackage{parskip}
\usepackage{amsmath,amssymb}
\usepackage{graphicx}
\usepackage{hyperref}
\usepackage{booktabs}
\usepackage{enumitem}
\usepackage{cite}

% ── Draft reference links ────────────────────────────────────────────────────
% Each citation carries a small superscript marker that opens the local copy of
% the paper in paper/refs/<bibkey>.pdf. Keys with no local PDF are left
% untouched, so the marker doubles as an inventory of what we hold.
% These are working notes, so the links stay on. Set \draftlinksfalse to strip.
\newif\ifdraftlinks
\draftlinkstrue

\ifdraftlinks
  \usepackage{etoolbox}
  \usepackage{xcolor}
  \newcommand{\refpdfdir}{refs}
  \newcommand{\refpdfscheme}{run:}
  \newcommand{\refpdfglyph}{\textbullet}
  \newcommand{\dlmark}[1]{%
    \IfFileExists{\refpdfdir/#1.pdf}%
      {\href{\refpdfscheme\refpdfdir/#1.pdf}%
            {\textsuperscript{\tiny\textcolor{blue}{\refpdfglyph}}}}%
      {}%
  }
  \let\dlorigcite\cite
  \renewcommand{\cite}[2][]{%
    \ifstrempty{#1}{\dlorigcite{#2}}{\dlorigcite[#1]{#2}}%
    \forcsvlist{\dlmark}{#2}%
  }
\fi

% Online locators. \detokenize keeps underscores literal, so DOIs such as
% 10.1007/978-3-319-55372-6_19 need no hand-escaping.
\newcommand{\doilink}[1]{%
  \href{https://doi.org/\detokenize{#1}}{\texttt{\footnotesize doi:\detokenize{#1}}}}
\newcommand{\arxivlink}[1]{%
  \href{https://arxiv.org/abs/\detokenize{#1}}{\texttt{\footnotesize arXiv:\detokenize{#1}}}}
\newcommand{\weblink}[2]{\href{#1}{\texttt{\footnotesize #2}}}
\newcommand{\nolink}[1]{{\footnotesize\itshape #1}}

\newlist{outline}{itemize}{4}
\setlist[outline]{leftmargin=1.2em,itemsep=2pt,parsep=1pt,topsep=3pt,label=--}
\newcommand{\Point}{\textbf{Point.}\ }
\newcommand{\Open}{\textbf{Open.}\ }

\sloppy
\hypersetup{colorlinks=true,linkcolor=blue!50!black,citecolor=blue!50!black,
            urlcolor=blue!50!black}

\title{Classical Modelling Notes\\
  \large Levels 1 and 2 for the OCEANS 2026 USV Identification Paper}
\author{Brian S. Bingham}
\date{\today}

\begin{document}
\maketitle

\begin{center}
\fbox{\parbox{0.9\textwidth}{\small
\textbf{Internal working notes. Not for circulation.}\\[0.4em]
This document holds the detail: the annotated reference checklist, the
development of the sufficiency argument and the derivation behind the
margin-based diagnostic. It is the scratch space for the traditional
engineering models, Levels 1 and 2.\\[0.4em]
The companion \texttt{paper.tex} is the team-facing document. It carries the
paper outline at outline-level specificity, plus section status and action
assignments. Detail belongs here; coordination belongs there. When something
here settles, promote a one-line summary into \texttt{paper.tex} rather than
copying the argument across.\\[0.4em]
Blue superscript dots mark citations with a local PDF in \texttt{refs/}.}}
\end{center}

\tableofcontents
\bigskip

\section{Introduction}
\label{sec:intro}

\begin{outline}
\item \textbf{Motivate.} Small USVs (under 7~m, single thruster plus rudder) are
  increasingly deployed for survey, inspection and autonomy research. Low-level
  speed and heading control is a prerequisite for any higher-level autonomy.
\item \textbf{The tension.} Control design requires a model, but the ``right''
  model depends on what the model is for. ``How much fidelity is enough?'' has no
  answer until ``enough for what?'' is specified.
  \begin{outline}
  \item Because that specification is usually left implicit, the default is a bias
    toward more fidelity: the benefit of an added term is easy to measure
    ($\Delta R^2$), while its cost is diffuse and rarely accounted for, namely
    identifiability from feasible experiments, experimental burden and loss of
    design insight.
  \item \textbf{Broadly, this is the overfitting problem.} Driving $R^2$ up is a naive
    objective once it is separated from what the model is for. Past some point the
    added terms improve the fit and damage the thing we actually want, which is good
    control. Machine learning at least made that trade visible by convention, through
    held-out data, and pays it down through
    regularization~\cite{suarez2026regularized}.
  \item Held-out scoring is a weaker check here than it is in machine learning,
    because the \emph{parameters} and not only the predictions are the deliverable. A
    maneuvering model can pass a prediction test while its individual coefficients are
    meaningless~\cite{alexandersson2022system}. Hence diagnostics that compare
    \emph{across} model structures rather than scoring one structure on held-out data.
  \item Prior work is \emph{not} naive about this, and the most careful example shows
    exactly where the difficulty sits. Sonnenburg and
    Woolsey~\cite{sonnenburg2013modeling} build genuine quantitative machinery for
    sufficiency: a cost function evaluated on held-out validation maneuvers,
    tabulated per model and per regime. The final decision is nonetheless made by
    judgment, discarding the model that scores best because its two extra parameters
    are ``not worth the slight edge'' in turn-rate estimation.
  \item \textbf{This is not a lapse on their part.} Their cost function scores the
    \emph{benefit} of added parameters. No comparable metric exists for the
    \emph{cost}, so there is nothing for the judgment to be replaced by. The
    quantitative apparatus stops one step short of the decision it was built to
    support. That vacancy, not any author's rigor, is what this paper is aimed at.
    \emph{Detail in Sec.~\ref{sec:rw-usv}.}
  \item The residual gap is that sufficiency is resolved by \emph{selecting} one
    model, after which the discarded models stop contributing.
  \end{outline}
\item \textbf{Framing.} George Box - all models are wrong
  but some are useful~\cite{box1987empirical}: a spectrum of models, each wrong in
  different ways, is collectively more useful than any one complex model.
\item \textbf{Sufficient for what? The decisions a control-oriented model must
  support.} Models (and simulations) are means to an end, not ends in themselves. The
  end here is specific: transferring low-level stabilization control from simulation
  to the vehicle.
  \begin{outline}
  \item \textbf{The naive premise, stated so we can reject it.} The implicit default is often to assume that the model is faithful replica of the system under test, a
    twin (analog or digital), so that with sufficient fidelity
    one could engineer control and autonomy entirely in simulation and transfer
    zero-shot to hardware. Fidelity is then an unbounded good and the only question
    is how much of it you can afford.
  \item \textbf{Fidelity is not the objective, and pursuing it past a point does
    concrete harm.} What we want is control that works on the vehicle, and that is
    not monotone in model fidelity. Four costs, none hypothetical:
    \begin{outline}
    \item \emph{The estimates get worse, not merely more numerous.} Added terms
      degrade the conditioning of the fit, so parameters that were previously well
      determined become less certain~\cite{suarez2026regularized}. That is a loss
      taken, not a gain forgone.
    \item \emph{Field time is the scarce resource.} Every added term needs excitation
      to identify it, so fidelity converts directly into trials on a platform where
      trials are the binding constraint.
    \item \emph{Design insight degrades.} A model with many terms predicts without
      explaining. It cannot say which term drives the behavior, which is exactly what
      architecture selection (D1) needs from it.
    \item \emph{False confidence.} A high-fidelity model invites the belief that the
      computed gains will transfer, which is D5 below and does not hold. Higher
      fidelity can therefore make a team \emph{less} likely to plan field validation.
      That is a schedule and safety risk, not a modeling one.
    \end{outline}
  \item \textbf{Two independent bounds on useful fidelity.} The economic bound below
    is the obvious one. The sharper claim is that even if fidelity were free there
    would still be a ceiling: once the model discrepancy falls inside the robustness
    margin of the design, further fidelity changes no decision and is inert. The
    limit is structural, not budgetary. \emph{This is the quantity Diagnostic~4
    (Sec.~\ref{sec:diagnostics}) computes.}
  \item \textbf{Two distinct claims, and the second is the one we need.}
    \begin{outline}
    \item \emph{Box's claim is epistemic.} Correctness is unattainable, so a wrong
      model cannot be made ``\,`correct' \dots by excessive elaboration,'' and
      ``overelaboration and overparameterization is often the mark of
      mediocrity''~\cite{box1976science}. Parsimony follows from the impossibility of
      truth.
    \item \emph{Right versus effective is a decision claim, and stronger.} Additional
      correctness is often genuinely available and should still be declined when it
      trades against the outcome. Skelton put this in control terms in
      1989~\cite{skelton1989model}: model errors ``need not be `small' but simply
      `appropriate' '' for the control design. Fidelity is not the criterion, fitness
      for the design decision is.
    \item Diagnostic~4 sits squarely on the second claim. Level~2 genuinely is a
      better description of the physics than Level~1, and past the margin threshold
      it still changes no design decision. Box does not predict that. Skelton does,
      and neither of them says how to check it on a given vehicle.
    \end{outline}
  \item \textbf{The gap we address.} Skelton's criterion has been available for
    thirty-five years and remains qualitative, since ``appropriate'' is not a
    computable predicate. The identification-for-control
    tradition~\cite{ljung1999system,gevers2005identification} formalized the idea for
    experiment design without producing a field test a small-boat practitioner can
    run. Diagnostic~4 is an attempt at that test for this vehicle class.
  \item \textbf{What sets how much model you need is the economics, not the physics:}
    the cost of the platform, access to the test environment and above all the ratio
    of hardware-test cost to simulation-test cost. A cheap vehicle on an accessible
    lake and an expensive vehicle with a narrow at-sea window warrant different
    answers to an identical modeling question.
  \item \textbf{What a model of this class \emph{can} support}, in rough order of
    increasing demand:
    \begin{outline}
    \item \textbf{D1: choose a control architecture.} Weakest requirement. The model
      need only get the dominant structure right: how many states, which couplings
      matter, where the dominant timescale sits.
    \item \textbf{D2: verify the architecture end-to-end.} Does the control system
      function, and function correctly, when wired up whole? Catches structural and
      integration faults on the bench rather than on the water.
    \item \textbf{D3: compare architectures relatively.} What is the cost and benefit
      of adding a feedforward term or a rate filter? Differences between
      architectures survive model error far better than absolute values do. Caveat:
      each candidate must be tuned before it can be compared, so tuning quality
      confounds the comparison unless explicitly controlled for.
    \item \textbf{D4: rehearse the tuning process.} Practice the procedure and build
      the tooling that will be used to tune in the field. The deliverable is the
      tuning workflow, not the gains it produces.
    \end{outline}
  \item \textbf{What the classical side contributes back.} Gain and phase margin state
    how much model error a design tolerates. That is a cost-side metric of exactly the
    kind missing from the sufficiency literature above: once the discrepancy between
    two model levels falls inside the margin of the controller designed from the
    simpler one, the extra fidelity buys nothing for that design. The margin is the
    transfer budget. \emph{Developed as Diagnostic~4, Sec.~\ref{sec:diagnostics}.}
  \item \textbf{What it \emph{cannot} support (D5):} identifying parameters and
    designing the controller in simulation and then expecting zero- or one-shot
    transfer. Saying this plainly is unusual, and it is the honest answer to ``why
    model at all if you re-tune in the field anyway.'' You are buying D1--D4. You are
    not buying the gains.
  \item Our three levels map onto these decisions rather than onto a fidelity ladder:
    Level~1 serves D1, Level~2 serves D2 and D3, Level~3 bounds how far D4 could be
    pushed with more data. \emph{Returned to in Sec.~\ref{sec:discussion}.}
  \item \Open This ladder merges the earlier purpose framing (choose / compare /
    parameterize) with the sim-to-real framing. D4 is the piece with no counterpart
    in the earlier version, and is arguably the most defensible claim we can make,
    since it is the one we can demonstrate rather than assert.
  \item \Open Placement. This is long for an introduction bullet. Candidate: promote to a short standalone section before
    Sec.~\ref{sec:methodology}, so the methodology can be presented as answering
    D1--D4 explicitly.
  \item \Open Scope of the claim. Is the economics argument (hardware-test cost vs.\
    simulation-test cost) something we can support in this paper, or a position
    statement we should flag as future work? It currently reads as the former.
  \end{outline}
\item \textbf{Scope.} Single thruster plus rudder, under 7~m, field-identified
  parameters.
  \begin{outline}
  \item Underactuated rudder steering is the more interesting case: control
    authority is an explicit function of forward speed, so the steering and speed
    channels cannot be treated as independent design problems. Differential or
    holonomic thrust reframes the problem rather than simplifying it, and removes
    precisely this coupling.
  \item It is also the more representative case for the installed base of small
    craft, even though the USV literature is dominated by differential-thrust
    research platforms.
  \item Distinct as well from large vessels, where scale-model tow-tank testing is a
    reasonable route to parameters, so identification is not the binding constraint.
  \end{outline}
\item \textbf{Contributions.}
  \begin{outline}
  \item A three-level model spectrum methodology for this class of USV.
  \item Cross-model diagnostics that identify structurally salient dynamics.
  \item A margin-based sufficiency test (Diagnostic~4) that reports the operating
    point at which added model fidelity stops changing the control design, replacing
    a judgment call with a computed number.
  \item An empirical demonstration of when physical structure is required (black-box
    simulation failure under limited data).
  \item A targeted data collection framework motivated by diagnostic gaps.
  \end{outline}
\end{outline}

% ─────────────────────────────────────────────────────────────────────────────
\section{Closely Related Work}
\label{sec:related}

Organizing axis: what is the least complex model structure that still supports
control design, and what does the available excitation actually permit one to
identify? Together these bound the cost of added fidelity, which is the quantity
that is missing when complexity is chosen by fit alone.

\subsection{Response Models as the Minimal Control-Oriented Description}
\label{sec:rw-response}

\begin{outline}

\item \textbf{Nomoto, Taguchi, Honda, Hirano (1957)}~\cite{nomoto1957steering},
  \emph{On the Steering Qualities of Ships}, Int.\ Shipbuilding Progress 4(35),
  354--370. \nolink{No digital copy in circulation.}
  \begin{outline}
  \item \Point Origin of the first-order steering model
    $r(s)/\delta(s)=K/(Ts+1)$. Collapses coupled sway-yaw into a gain and a time
    constant that map directly onto compensator decisions: $T$ sets achievable
    closed-loop bandwidth, $K$ sets loop gain. This is exactly our Level-1
    structure, applied to surge as well as yaw.
  \item \Open Developed for large displacement vessels. Does the first-order form
    hold for a sub-7~m hull? Sonnenburg's data say yes at higher speeds; we should
    state this rather than assume it.
  \end{outline}

\item \textbf{Fossen (2011)}~\cite{fossen2011handbook},
  \emph{Handbook of Marine
  Craft Hydrodynamics and Motion Control}, Wiley. \doilink{10.1002/9781119994138}
  \begin{outline}
  \item \Point Canonical reduction of the linearized maneuvering equations to the
    Nomoto form and the source of our notation. Situates the response model in the
    broader hierarchy, which is the hierarchy we are traversing.
  \item \Open We hold the 2011 first edition; a 2021 second edition exists
    (\doilink{10.1002/9781119575016}). Cite whichever we actually work from.
  \end{outline}

\item \textbf{{\AA}str{\"o}m and K{\"a}llstr{\"o}m (1976)}~\cite{astrom1976identification},
  \emph{Identification of Ship Steering Dynamics}, Automatica 12(1), 9--22.
  \doilink{10.1016/0005-1098(76)90064-9}
  \begin{outline}
  \item \Point Founds the identification side of this tradition: linear steering
    dynamics estimated from full-scale trial records by maximum likelihood.
    Supports our central claim that the binding constraint is experimental, not
    algorithmic.
  \item \Open They use ML with a Kalman filter; we use ordinary least squares on a
    discrete ARX(1,1). Do we need to justify OLS, or is it self-evident at 10~Hz on
    a short clean record? Probably one sentence.
  \end{outline}

\item \textbf{K{\"a}llstr{\"o}m and {\AA}str{\"o}m (1981)}~\cite{kallstrom1981experiences},
  \emph{Experiences of System Identification Applied to Ship Steering},
  Automatica 17(1), 187--198. \doilink{10.1016/0005-1098(81)90094-7}
  \nolink{No local copy. Try Lund LUP; the 1976 companion is free there.}
  \begin{outline}
  \item \Point Practical follow-on: the maneuvers a vessel can safely execute limit
    which parameters are distinguishable. This is the cleanest available statement
    of the excitation-limits-identifiability point.
  \item \Open Need the PDF to confirm this claim is actually made in these terms
    before we lean on it. Currently sourced from secondary description.
  \end{outline}

\end{outline}

\subsection{Physics-Based Maneuvering Models and the Identifiability Barrier}
\label{sec:rw-physics}

\begin{outline}

\item \textbf{Abkowitz (1964)}~\cite{abkowitz1964lectures},
  \emph{Lectures on
  Ship Hydrodynamics: Steering and Manoeuvrability}, Report Hy-5, Lyngby.
  \nolink{No digital copy in circulation.}
  \begin{outline}
  \item \Point The canonical high-fidelity structure: third-order truncated Taylor
    expansion of hydrodynamic forces about steady straight-line motion. Cited for
    provenance as the pole opposite Nomoto.
  \item \Open Cite directly, or via a secondary account (Fossen Ch.~7, Sonnenburg
    diss.\ Sec.~2.1) that we can actually read? Leaning secondary.
  \end{outline}

\item \textbf{Yasukawa and Yoshimura (2015)}~\cite{yasukawa2015mmg},
  \emph{Introduction of MMG Standard Method for Ship Maneuvering Predictions},
  J.\ Marine Sci.\ Tech.\ 20(1), 37--52. \doilink{10.1007/s00773-014-0293-y}
  \nolink{Open access; Springer redirect blocked scripted retrieval.}
  \begin{outline}
  \item \Point The modular alternative to Abkowitz and the current standard for
    ship maneuvering prediction. Establishes that the high-fidelity pole is a live
    practice, not a historical artifact.
  \item \Open Does MMG earn its space in a 6-page paper, or is Abkowitz alone
    enough to anchor the pole? Candidate cut if we run long.
  \end{outline}

\item \textbf{Su{\'a}rez, Berndt, Abdel-Maksoud (2026)}~\cite{suarez2026regularized},
  \emph{Regularized Machine Learning for System Identification of Ship
  Free-Running Manoeuvres from CFD-Based Synthetic Data}. \arxivlink{2606.17121}
  \begin{outline}
  \item \Point \emph{The citation our cost-of-fidelity argument rests on.}
    Abkowitz derivatives are strongly multicollinear; conditioning degrades as the
    coefficient set grows; the effective remedies are regularization and more
    diverse excitation, not better optimization. ``Conditioning degrades as the
    coefficient set grows'' is the cost of added fidelity stated as a measurable
    quantity.
  \item \Point Their data are CFD-generated from idealized zigzag and turning
    maneuvers, which sharpens rather than weakens the point: if a coefficient set is
    ill-conditioned under noise-free synthetic excitation designed expressly for the
    task, it will not be recovered from a short trial on the water.
  \item \Open Preprint, under review. If it has not landed by submission, restore
    Luo and Li (2017), \doilink{10.1016/j.apor.2017.06.008}, alongside it rather
    than instead of it; that paper attributes parameter drift to dynamic
    cancellation among linear derivatives plus multicollinearity among nonlinear
    ones.
  \item \Open Do we run Ridge on our own Level-2 parameters as a further diagnostic?
    Cheap to do, and it would make the citation reciprocal. Tracked at
    Sec.~\ref{sec:l2}; note that Diagnostic~4 is already taken by the margin test.
  \end{outline}

\item \textbf{Xu and Guedes Soares (2025)}~\cite{xu2025review},
  \emph{Review of
  System Identification for Manoeuvring Modelling of Marine Surface Ships},
  J.\ Marine Sci.\ Appl.\ 24, 459--478. \doilink{10.1007/s11804-025-00681-w}
  \begin{outline}
  \item \Point Single citation covering ``this difficulty is general, not specific
    to one method or one hull.'' Saves us three or four individual citations.
  \item \Open Confirm it addresses small USVs and not only ship-scale vessels. If
    ship-only, we must say so rather than imply generality we have not checked.
  \end{outline}

\item \textbf{Alexandersson, Mao, Ringsberg (2022)}~\cite{alexandersson2022system},
  \emph{System Identification of Vessel Manoeuvring Models}, Ocean Engineering
  266, 112940. \doilink{10.1016/j.oceaneng.2022.112940}
  \begin{outline}
  \item \Point Identified coefficients need not be individually meaningful even
    when aggregate prediction is good. This separates ``predicts well'' from
    ``parameters mean something,'' which is precisely the interpretability axis our
    three levels are arranged along.
  \item \Open Confirm the vessel scale and data source (model tests vs.\ full
    scale) so we characterize the result accurately.
  \end{outline}

\end{outline}

\subsection{Control-Oriented Identification of Small Uncrewed Craft}
\label{sec:rw-usv}

\begin{outline}

\item \textbf{Sonnenburg and Woolsey (2013)}~\cite{sonnenburg2013modeling},
  \emph{Modeling, Identification, and Control of an Unmanned Surface Vehicle},
  J.\ Field Robotics 30(3), 371--398. \doilink{10.1002/rob.21452}
  \\ Companion: \textbf{Sonnenburg \emph{et al.}\ (2010)}~\cite{sonnenburg2010control},
  OCEANS 2010 Seattle. \doilink{10.1109/OCEANS.2010.5664297}
  \nolink{Local copies are the 2013 dissertation (superset) and the extended VT report.}
  \begin{outline}
  \item \Point \emph{The nearest antecedent.} A 4.8~m planing RHIB; objective stated
    as a model ``sufficiently rich for model-based control design, sufficiently
    simple to allow parameter identification.'' Candidate speed and steering models
    identified and compared from field trials.
  \item \Point Two-stage identification: steady-state input-output maps from
    constant-input maneuvers first, then transient parameters with those gains
    fixed. Our two-phase separation of quadratic drag from thrust gain follows the
    same logic of estimating each parameter where it dominates. \emph{Cite again in
    Sec.~\ref{sec:l2}.}
  \item \Point Their conclusion is that no single model wins: Nomoto-with-sideslip
    below roughly 2.6~m/s, plain Nomoto above, deployed as a speed-scheduled
    combination.
  \item \Point \emph{How they evaluate sufficiency} (answers the question raised in
    Sec.~\ref{sec:intro}). Candidates are scored by a cost function evaluated on
    held-out validation maneuvers and tabulated per model and per regime, for both
    forward and sternward motion. Selection then applies a parsimony judgment: the
    full linear steering model wins marginally on cost in the stable regimes, but its
    two extra parameters are judged ``not worth the slight edge in turn rate
    estimation.''
  \item \Point \emph{This is the point we build on.} Their cost function measures the
    benefit of added parameters. Nothing in the procedure measures the cost, and no
    accepted metric for it exists, so the trade is necessarily settled by expert
    intuition. Read generously, this is the strongest prior work reaching the limit
    of the available apparatus rather than declining to use it. Our cross-model
    diagnostics are an attempt to put evidence on the vacant side of that scale.
  \item \Point Two differences define our contribution. Their vessel steers by
    vectored thrust, which removes the speed-dependent control effectiveness that
    dominates rudder-steered craft. And their comparison stops at structured
    models, with no data-driven level to bound residual complexity.
  \item \Open Their envelope reaches 22~kn with regime transitions; ours is 0--3~m/s
    in displacement only. Does their speed-scheduling conclusion transfer, or is it
    an artifact of crossing into planing?
  \item \Open Should we adopt their steady-state-first procedure for Level~1 too,
    rather than fitting ARX directly? Would make the two levels methodologically
    consistent.
  \end{outline}

\item \textbf{Eriksen and Breivik (2017)}~\cite{eriksen2017modeling},
  \emph{Modeling, Identification and Control of High-Speed ASVs: Theory and
  Experiments}, LNCIS 474, 407--431. \doilink{10.1007/978-3-319-55372-6_19}
  \begin{outline}
  \item \Point Demonstrates the full path from field identification to deployed
    control: an 8.45~m ASV spanning displacement, semi-displacement, and planing,
    with four controllers tested in full scale. Evidence that control-oriented
    identification pays off, which is the premise we inherit.
  \item \Open At 8.45~m with rich instrumentation. Which parts of the method
    survive our sensor suite (GPS speed at 5~Hz, IMU yaw rate)? Worth one explicit
    sentence rather than silence.
  \end{outline}

\item \textbf{Muske, Ashrafiuon, Haas, McCloskey, Flynn (2008)}~\cite{muske2008identification},
  \emph{Identification of a Control Oriented Nonlinear Dynamic USV Model},
  ACC 2008, 562--567. \doilink{10.1109/ACC.2008.4586551}
  \begin{outline}
  \item \Point \emph{Direct precedent for our Level-2 method.} Drag parameters
    identified from dedicated deceleration and terminal-velocity tests, with thrust
    deliberately excluded from the regression. Same reasoning as our coasting-segment
    identification of $d_2$. \emph{Cite again in Sec.~\ref{sec:l2}.}
  \item \Open They use a model boat in controlled conditions. Does coast-down
    survive wind and current in the field? Our dataset has limited coasting; is one
    long deceleration enough, and how do we defend that?
  \end{outline}

\item \textbf{Hann, Sirisena, Wongvanich (2010)}~\cite{hann2010simplified},
  \emph{Simplified Modeling Approach to System Identification of Nonlinear Boat
  Dynamics}, ACC 2010, 5218--5223. \doilink{10.1109/ACC.2010.5530459}
  \nolink{No local copy.}
  \begin{outline}
  \item \Point Simplified identification of nonlinear boat dynamics with
    deliberately tractable experiments.
  \item \Open Unclear how much this adds beyond Muske. Candidate cut if the
    overlap is large; decide once we have read it.
  \end{outline}

\item \textbf{Wirtensohn, Reuter, Blaich, Schuster, Hamburger (2013)}~\cite{wirtensohn2013modelling},
  \emph{Modelling and Identification of a Twin Hull-Based Autonomous Surface
  Craft}, MMAR 2013, 121--126. \doilink{10.1109/MMAR.2013.6669892}
  \begin{outline}
  \item \Point Concrete example of a complete identification campaign at small-ASV
    scale. Useful as the template we are deviating from.
  \item \Open Differential thrust, so control effectiveness is not speed-dependent
    the way ours is. This supports our scope claim in Sec.~\ref{sec:rw-position};
    make sure we use it there rather than overstating its similarity here.
  \end{outline}

\item \textbf{Morel, Orihuela, Bejarano (2023)}~\cite{morel2023modelling},
  \emph{Modelling and Identification of an Autonomous Surface Vehicle}, MARTECH
  2023, Castell{\'o} de la Plana.
  \\ Journal version: \textbf{Morel, Orihuela, Combastel, Bejarano
  (2025)}~\cite{morel2025practical}, \emph{Practical Identification Approach for the
  Actuation Dynamics of Autonomous Surface Vehicles with Minimal Instrumentation},
  Ocean Engineering 319, 120098. \doilink{10.1016/j.oceaneng.2024.120098} /
  \arxivlink{2410.00631}
  \begin{outline}
  \item \Point \emph{Closest match to our instrumentation constraint.} Premise is
    that the sensor suite of a small autonomous craft, not the estimator, is the
    binding constraint. The 2023 workshop paper names the vessel (Yellowfish ASV)
    and tabulates parameters; the 2025 journal paper generalizes the method.
  \item \Open They model \emph{actuation} dynamics explicitly. Do we need an
    actuator model, or is our normalized RCOU-to-force mapping adequate? If we skip
    it, say why, because this paper makes the omission conspicuous.
  \end{outline}

\item \textbf{Tanveer and Ahmad (2022)}~\cite{tanveer2022yaw},
  \emph{Unmanned
  Surface Vehicle: Yaw Modeling and Identification}, ICEANS 2022.
  \arxivlink{2303.13064}
  \begin{outline}
  \item \Point Honest illustration of what a single short experiment supports: a
    second-order yaw model from pool trials at 59\% cross-validation fit. Supports
    our data-coverage claim from the low end.
  \item \Open Twin thruster, pool trials, two-page paper. Is it too weak to carry
    weight, or is its weakness exactly the point? Currently used as the latter.
  \end{outline}

\item \textbf{Synthesis for this subsection.} The literature converges on
  simplified, field-identified models as the right instrument at this scale. What it
  does not do is \emph{retain} more than one. Each study selects a structure,
  identifies it, and proceeds to control; comparison, where it occurs, is used to
  discard alternatives rather than to interrogate the dynamics.

\end{outline}

\subsection{Data-Driven and Grey-Box Alternatives}
\label{sec:rw-datadriven}

\begin{outline}

\item \textbf{Woo, Park, Yu, Kim (2018)}~\cite{woo2018dynamic},
  \emph{Dynamic
  Model Identification of Unmanned Surface Vehicles Using Deep Learning Network},
  Applied Ocean Research 78, 123--133. \doilink{10.1016/j.apor.2018.06.011}
  \nolink{No local copy. Highest-priority acquisition.}
  \begin{outline}
  \item \Point LSTM reduces surge error 76.9\% and yaw-rate error 60.7\% versus
    linear models. Our Level-3 result directly challenges the generality of this
    headline, so the paper is the foil for our central empirical claim.
  \item \Open \emph{Must verify before we assert it:} do they report one-step-ahead
    prediction only, with no closed-loop simulation? Our Sec.~\ref{sec:rw-position}
    gap~2 depends entirely on this being true. Do not submit without checking.
  \end{outline}

\item \textbf{Xu, Han, Cheng, Cheng, Ge (2022)}~\cite{xu2022pinn},
  \emph{A
  Physics-Informed Neural Network for the Prediction of Unmanned Surface Vehicle
  Dynamics}, JMSE 10(2), 148. \doilink{10.3390/jmse10020148}
  \begin{outline}
  \item \Point Embedding 3-DOF structure in the training loss improves
    generalization from small datasets. Supports our ``structure is regularization''
    claim from the opposite direction: they add structure to a network, we measure
    what its absence costs.
  \item \Open How small is their ``small''? If it is far larger than our 130~s, the
    comparison needs qualifying.
  \end{outline}

\item \textbf{Zhang, Li, Xiong, He (2024)}~\cite{zhang2024gbm},
  \emph{GBM-ILM:
  Grey-Box Modeling Based on Incremental Learning and Mechanism for Unmanned Surface
  Vehicles}, JMSE 12(4), 627. \doilink{10.3390/jmse12040627}
  \begin{outline}
  \item \Point Represents the current hybrid state of the art: mechanism model plus
    incremental learning. Establishes that the field's answer to data hunger is to
    reintroduce structure, which is our answer too, arrived at by a different route.
  \item \Open Is this the best representative, or is there a stronger 2025--26
    grey-box USV paper? One search before submission.
  \end{outline}

\item \textbf{Xue, Wang, Qin, Deng (2022)}~\cite{xue2022probabilistic},
  \emph{Probabilistic Identification of Unmanned Surface Vehicles Using Efficient
  Gaussian Processes with Uncertainty Propagation}, ICUS 2022, 1004--1009.
  \doilink{10.1109/ICUS55513.2022.9986638} \nolink{No local copy.}
  \begin{outline}
  \item \Point An alternative answer to ``how much should you trust this model'':
    propagate uncertainty so model confidence is available to the controller.
  \item \Open Lowest-priority citation in the section. Cut if we run long.
  \end{outline}

\item \textbf{Synthesis for this subsection.} Unstructured learning is data-hungry
  in exactly the way field identification of a small USV cannot accommodate. Our
  Level-3 NARX is not a competitor to these methods but a deliberately
  unregularized reference point, whose closed-loop failure quantifies how much
  regularization the physical structure supplies.

\end{outline}

\subsection{Positioning of This Work}
\label{sec:rw-position}

\begin{outline}
\item \textbf{Gap 1: hierarchy as selection, not retention.} Prior work retains one
  level and discards the rest. We retain all three and use disagreement between them
  as a diagnostic that localizes which dynamics are structurally significant and
  which are simply unobserved. Against~\cite{sonnenburg2013modeling,eriksen2017modeling}.
  \begin{outline}
  \item \Open Is this alone strong enough for an OCEANS contribution, or does it
    read as methodology-only? Gaps 2 and 3 may be doing more work.
  \end{outline}
\item \textbf{Gap 2: evaluation protocol.} Data-driven marine models are
  predominantly evaluated on one-step prediction, obscuring the closed-loop
  divergence that determines usefulness for control design. We report both metrics
  for the same model and dataset. Against~\cite{woo2018dynamic}.
  \begin{outline}
  \item \Open Contingent on the Woo verification above.
  \end{outline}
\item \textbf{Gap 3: vessel class.} The small, single-thruster, rudder-steered
  vessel is underrepresented: closest prior work uses vectored
  thrust~\cite{sonnenburg2013modeling}, differential
  thrust~\cite{wirtensohn2013modelling,tanveer2022yaw}, or a substantially larger
  platform~\cite{eriksen2017modeling}. Rudder steering makes control authority an
  explicit function of forward speed, the coupling our diagnostics target.
  \begin{outline}
  \item \Open Do we need a citation establishing that rudder-steered small USVs
    matter operationally, or is the platform's existence sufficient warrant?
  \end{outline}
\item \textbf{Framing.} The intended use is not only compensator synthesis but the
  simulation environments in which autonomy is developed and
  tested~\cite{bingham2019gazebo}. This aligns the work with identification for
  control~\cite{skelton1989model,ljung1999system,gevers2005identification}, where
  model quality is judged against the design decision the model informs rather than
  fit statistics alone.
  \begin{outline}
  \item \emph{Resolved:} the sim-to-real framing in Sec.~\ref{sec:intro} makes the
    simulator the primary consumer of these models, not a secondary one, so the
    citation is on-topic rather than incidental. Keep.
  \end{outline}
\end{outline}

% ─────────────────────────────────────────────────────────────────────────────
\section{Model Spectrum Methodology}
\label{sec:methodology}

\begin{outline}
\item \textbf{Framing.} Three levels, each identified from the same dataset. Each
  answers a different design question.
\end{outline}

\subsection{Level~1: First-Order Linear SISO Models}
\label{sec:l1}
\begin{outline}
\item Nomoto-style first-order transfer functions for surge and yaw, treated as
  independent SISO systems.
  \begin{outline}
  \item Surge: $G_v(s) = K/(\tau s + 1)$, input throttle (normalized), output speed.
  \item Yaw: $G_r(s) = K_r/(\tau_r s + 1)$, input rudder, output yaw rate.
  \end{outline}
\item \textbf{Identification.} Discrete-time ARX(1,1) by ordinary least squares;
  convert poles to physical $\tau$ and $K$.
\item \textbf{Design utility.} Directly readable time constants and gains for PID
  tuning; reveals dominant timescale per channel.
\item \textbf{Limitation.} Linear, decoupled; ignores quadratic drag, surge-yaw
  coupling, and operating-regime dependence. Alone, risks a compensator
  architecture mismatch.
\item \Open Adopt Sonnenburg's steady-state-first procedure here for consistency
  with Level~2? See Sec.~\ref{sec:rw-usv}.
\end{outline}

\subsection{Level~2: Nonlinear State-Space Model}
\label{sec:l2}
\begin{outline}
\item Physics-based two-state model from planar rigid-body equations:
  \begin{outline}
  \item Surge: $\dot{v} = k_T u - d_2 v|v|$.
  \item Yaw: $\dot{r} = k_\delta \delta - d_r r$.
  \end{outline}
\item Quadratic drag $v|v|$ is the key nonlinearity relative to Level~1.
\item \textbf{Two-phase identification} (precedent:
  \cite{muske2008identification,sonnenburg2013modeling}):
  \begin{outline}
  \item Phase 1 (coasting, $u\approx 0$): identify $d_2$ from zero-thrust
    deceleration segments.
  \item Phase 2 (powered): identify $k_T$ with $d_2$ fixed.
  \item Decoupling avoids conflating thrust and drag in a joint fit.
  \end{outline}
\item \textbf{Design utility.} Reveals the drag structure relevant to feedforward
  compensation and speed-hold steady state, while retaining physical
  interpretability.
\item \Open Report a conditioning or regularization check on the joint fit, per
  \cite{suarez2026regularized}, to make the identifiability argument concrete on our
  own data rather than only by citation.
\end{outline}

\subsection{Level~3: Black-Box NARX Model}
\label{sec:l3}
\begin{outline}
\item NARX multilayer perceptron predicting $[v(k{+}1), r(k{+}1)]$ from $n$ lags of
  $[v, r, \text{throttle}, \text{rudder}]$. MIMO, so it can discover coupling the
  structured models assume away.
\item \textbf{Architecture.} 3 lags, one hidden layer (20 neurons), ReLU,
  StandardScaler, early stopping.
\item \textbf{Training.} Chronological 80/20 split (first 100~s train, last 30~s
  test) to avoid temporal leakage.
\item \textbf{Evaluation.} One-step prediction \emph{and} closed-loop simulation
  reported separately; these diverge significantly at this data scale.
\item \textbf{Design utility.} Bounds the benefit of added complexity given
  available data; motivates richer collection; reference point as the dataset grows.
\end{outline}

\subsection{Cross-Model Diagnostics}
\label{sec:diagnostics}
\begin{outline}
\item The spectrum's value is the comparative analysis across levels, not the three
  fits individually.
\item \textbf{Diagnostic 1: speed-dependent rudder gain.} Does replacing constant
  $K_r$ with $K_r v$ improve the yaw-rate fit? Result: no improvement with current
  data ($\Delta R^2 = -0.03$), which flags a gap in data coverage rather than
  absence of the effect.
\item \textbf{Diagnostic 2: drag nonlinearity.} Linear vs.\ quadratic fit to
  coasting deceleration. Result: quadratic dominant, $\Delta R^2 = +0.20$,
  $D_1 \approx 0$, $D_2 = 0.44$~m$^{-1}$.
\item \textbf{Diagnostic 3: rudder-induced surge drag.} Correlation of $\delta^2$
  with surge residuals. Result: negligible, $r = -0.08$; coupling term not
  warranted.
\item \textbf{Diagnostic 4: is the added fidelity inside the robustness margin?}
  The first three diagnostics ask whether a term is \emph{present} in the data. This
  one asks whether it \emph{changes the design}, which is the question
  Sec.~\ref{sec:intro} argues has no accepted metric. It does have one: the
  robustness margin of the compensator is a budget for model error, so the test is
  whether the Level-1 to Level-2 discrepancy fits inside it.
  \begin{outline}
  \item \textbf{Step 1.} Design a nominal compensator from Level~1 alone (PI on
    surge, P or PI on yaw). Record gain margin, phase margin and the complementary
    sensitivity $T(j\omega)$.
  \item \textbf{Step 2.} Linearize Level~2 about operating points spanning the tested
    envelope. For surge the quadratic drag gives a speed-scheduled first-order plant,
    $\tau(v_0) = 1/(2 d_2 v_0)$ and $K(v_0) = k_T/(2 d_2 v_0)$, so both the gain and
    the time constant fall as $1/v_0$. With the identified $d_2 = 0.438$ and
    $k_T = 3.963$ this is roughly $\tau: 1.1 \rightarrow 0.4$~s and
    $K: 4.5 \rightarrow 1.6$ across $v_0 = 1$ to $2.9$~m/s, against Level~1's fixed
    $\tau = 1.95$~s, $K = 6.42$.
  \item \textbf{Step 3.} Form the multiplicative error
    $\ell(\omega, v_0) = |G_2(j\omega, v_0)/G_1(j\omega) - 1|$ and test robust
    stability, $\|\ell T\|_\infty < 1$. The cheap version that will fit in a short
    paper is the gain ratio $K(v_0)/K_1$ against the gain margin.
  \item \textbf{Step 4, the useful output.} The test does not return yes or no. It
    returns \emph{the speed at which Level~1 stops being sufficient}, which is an
    actionable number for a control designer and a direct answer to ``enough for
    what.''
  \item \textbf{Expected findings, to be confirmed.} For yaw, Level~2 reduces to
    $r/\delta = 0.771/(0.523 s + 1)$ against Level~1's $0.771/(0.471 s + 1)$: the
    same DC gain and an 11\% difference in time constant, so $\ell$ is small and the
    test should pass across the envelope. Level~2 buys nothing for yaw control, and
    we should say so. For surge the gain falls by roughly a factor of four across the
    envelope, so the test should fail at the upper end. Note the failure mode is
    sluggishness rather than instability, because quadratic drag is stabilizing: the
    Level-1 design is conservative at speed, not unsafe.
  \item \Open Scope. The full $\|\ell T\|_\infty$ treatment may not fit this paper.
    Minimum viable version: report $K(v_0)/K_1$ against the Level-1 gain margin and
    name the crossover speed. Full treatment flagged as future work.
  \end{outline}
\item \textbf{Implication.} Diagnostics 1--3 inform the Level-2 structure and specify
  what the next data collection must target. Diagnostic~4 converts the sufficiency
  question from a judgment into an arithmetic one, and is the piece that lets this
  work claim a method rather than a position.
\end{outline}

% ─────────────────────────────────────────────────────────────────────────────
\section{Experimental Platform and Data Collection}
\label{sec:experiment}

\begin{outline}
\item \textbf{Vehicle.} [model/name], approx.\ [length]~m, single propeller plus
  rudder, ArduPilot autopilot, GPS and IMU.
  \begin{outline}
  \item Inputs: RC throttle (RCOU\_C3) and rudder (RCOU\_C1), 1100--1900~$\mu$s PWM.
  \item Outputs: GPS speed (5~Hz), IMU yaw rate (137~Hz).
  \end{outline}
\item \textbf{Data cleaning.} Multi-rate streams aligned to a 10~Hz common grid by
  linear interpolation; RC commands normalized to $[-1,1]$; one 29.3~s gap (vessel
  repositioning) handled by segmentation.
\item \textbf{Current dataset.} Single step-response trial, 130~s working window
  (150--280~s of a 341~s recording), throttle 0--39\%, speed 0--2.9~m/s.
  Intentionally minimal: a proof of concept for the workflow.
\item \textbf{Planned dataset (full paper).} Multiple trials spanning throttle
  0--100\%, dedicated rudder sweeps at fixed speeds to resolve speed-dependent gain,
  and maneuvers spanning displacement, semi-displacement, and planing.
  [describe vehicle speed range and transition speeds]
\end{outline}

% ─────────────────────────────────────────────────────────────────────────────
\section{Results}
\label{sec:results}

\subsection{Level~1: First-Order SISO}
\begin{outline}
\item Table: identified $\tau$ and $K$ for surge and yaw.
\item Simulation $R^2$ for both channels.
\item Figure: measured vs.\ simulated speed and yaw rate.
\item Note: parameters directly usable for PID design.
\end{outline}

\subsection{Level~2: Nonlinear State-Space}
\begin{outline}
\item Table: identified $k_T$, $d_2$, $k_\delta$, $d_r$.
\item Compare $d_2$ from coasting vs.\ joint fit; motivate why two-phase is
  necessary.
\item Simulation $R^2$ comparison vs.\ Level~1.
\item Figure: measured vs.\ simulated, combined plot across all three levels.
\end{outline}

\subsection{Level~3: NARX}
\begin{outline}
\item One-step prediction $R^2$ vs.\ closed-loop simulation $R^2$; highlight the gap.
\item Figure: simulation divergence for surge.
\item Interpretation: limited data makes black-box models unreliable simulators;
  physical structure is essential regularization at this data scale.
\end{outline}

\subsection{Cross-Model Diagnostic Summary}
\begin{outline}
\item Summary table: four diagnostics, findings, implication for model and control
  design.
\item Diagnostic~4 reported in its minimum viable form: $K(v_0)/K_1$ against the
  Level-1 gain margin, plus the crossover speed at which Level~1 stops being
  sufficient for surge.
\item Key message: the spectrum reveals structure that no single model exposes, and
  the margin test says which of that structure changes the design.
\end{outline}

% ─────────────────────────────────────────────────────────────────────────────
\section{Discussion}
\label{sec:discussion}

\begin{outline}
\item \textbf{What the spectrum reveals for this vehicle class.}
  \begin{outline}
  \item Surge and yaw are weakly coupled at the speeds tested, so decoupled SISO
    controllers are a reasonable starting point.
  \item Quadratic drag is significant: speed-hold controllers need nonlinear
    feedforward; linear controllers will show speed-dependent steady-state error.
  \item Speed-dependent rudder gain can neither be confirmed nor denied. This is the
    primary gap the expanded dataset must fill. If confirmed, gain scheduling or
    adaptive control is warranted.
  \end{outline}
\item \textbf{Implications for compensator architecture.}
  \begin{outline}
  \item Level~1 $\rightarrow$ baseline PID structure with identified gains.
  \item Level~2 $\rightarrow$ nonlinear feedforward for drag on the surge channel,
    where the margin test is expected to show Level~1 failing at speed. For yaw the
    two levels nearly coincide, so Level~2 changes nothing and we should say so.
  \item Gain scheduling on the rudder loop remains open, gated on Diagnostic~1 rather
    than on Level~2.
  \item Level~3 $\rightarrow$ motivates MPC or learning-based control once data
    suffices.
  \end{outline}
\item \textbf{Fidelity vs.\ insight: return to D1--D5 (Sec.~\ref{sec:intro}).}
  \begin{outline}
  \item State plainly that none of the three levels reaches D5. Even the
    highest-fidelity model here cannot eliminate experimental fine-tuning, and we
    should say so rather than let the reader infer otherwise.
  \item The defensible claim is D1 through D4: the spectrum tells you which
    architecture to build, lets you verify it end to end, supports relative
    comparison between candidates and tells you where in the parameter space field
    tuning will actually be needed so that the tuning workflow can be built in
    advance.
  \item \Open Is ``the model does not remove field tuning, but it tells you what to
    tune and lets you rehearse it'' a strong enough payoff to headline? It is honest,
    and it is the question most modeling papers decline to answer.
  \end{outline}
\item \textbf{Limitations} of the current dataset and what the full study adds.
\end{outline}

% ─────────────────────────────────────────────────────────────────────────────
\section{Conclusion}
\label{sec:conclusion}

\begin{outline}
\item Summarize the three-level spectrum methodology for small underactuated USVs.
\item Restate the core finding: the spectrum collectively provides more control
  design guidance than any single model; simple models are not a stepping stone but
  a persistent design tool.
\item State what the full dataset will allow: resolution of speed-dependent rudder
  dynamics, validation across operating regimes and a more definitive comparison of
  black-box vs.\ structured models.
\item \textbf{Future work: margin-based sufficiency.} Diagnostic~4 is reported here in
  its minimum viable form. The extension is to treat the robustness margin as the
  formal sufficiency criterion across the whole spectrum, computing for each design
  decision the model level at which further fidelity is provably inert. That would
  replace the expert judgment identified in Sec.~\ref{sec:rw-usv} with a computable
  bound, and it generalizes past this vehicle class.
\item Release the dataset and analysis notebooks as open tools for the USV research
  community.
\end{outline}

% ─────────────────────────────────────────────────────────────────────────────
\nocite{*}
\bibliographystyle{IEEEtran}
\bibliography{references}

\end{document}
